\renewcommand*{\authorOfNextLines}{Helmut Quinz}
Wesentlich für das Generieren des PDF-Dokumentes sind die Einstellungen für den Standardcompiler sowie das Standard-Bibliographieprogramm.
Als Compiler ist \gqq{\texttt{LuaLaTeX}} zu verwenden, für die Bibliographie \gqq{\texttt{Biber}}
\begin{figure}[ht]
	\begin{overpic}[width=130mm,abs,unit=1mm,,tics=10]{../13_D3/VorlageEinstellungenUebersetzen.png}
		\put(26,2)		{\opFrame				[white, rounded corners=1mm]{35mm,3mm}}
		\put(26,11.5)		{\opFrame				[white, rounded corners=1mm]{35mm,3mm}}
	\end{overpic}	
	\caption{\TeX -Editor Einstellungen für das Übersetzen}
\end{figure}
\FloatBarrier
Die Zeichenkodierung der Text-Dateien sollte in Grundeinstellung auf \texttt{utf8} gestellt sein. 
Es schadet jedoch nicht dies zu kontrollieren:
\begin{figure}[ht]
	\begin{overpic}[width=130mm,abs,unit=1mm,,tics=10]{../13_D3/VorlageEinstellungenEditor.png}
		\put(36,10)		{\opFrame				[white, rounded corners=1mm]{50mm,3mm}}
	\end{overpic}	
	\caption{\TeX -Editor Einstellungen für das Erstellen der Dateien}
\end{figure}
\FloatBarrier